\documentclass[11pt]{article}
\usepackage[margin=1in]{geometry}
\usepackage{amsmath, amssymb, amsthm}
\usepackage{mathtools}
\usepackage{enumitem}
\usepackage{hyperref}

\newtheorem{theorem}{Theorem}
\newtheorem{lemma}[theorem]{Lemma}
\newtheorem{proposition}[theorem]{Proposition}
\theoremstyle{definition}
\newtheorem{definition}{Definition}
\newtheorem{example}{Example}

\title{Relations and Orderings}
\author{Math Foundations --- Concept 03}
\date{}

\begin{document}
\maketitle

\section{Relations on a Set}

A \emph{binary relation} on a set $A$ is a subset $R \subseteq A \times A$.
We write $a \mathrel{R} b$ as shorthand for $(a,b) \in R$.

\subsection*{First, an example}
Before any abstract axioms, look at congruence mod $3$ on $A = \{0,1,2,3,4,5,6,7,8\}$. Declare $x \equiv y \pmod 3$ when $3 \mid (x-y)$. The equivalence classes are exactly the three residue buckets:
\begin{itemize}[leftmargin=2em]
  \item $[0] = \{0, 3, 6\}$
  \item $[1] = \{1, 4, 7\}$
  \item $[2] = \{2, 5, 8\}$
\end{itemize}
Notice three features: every element of $A$ lands in exactly one bucket (covering + disjoint $\Rightarrow$ a \emph{partition}); the buckets are determined by the relation, not chosen by hand; and the relation $\equiv \pmod 3$ is reflexive ($x - x = 0$), symmetric ($3 \mid (x-y) \Leftrightarrow 3 \mid (y-x)$), and transitive. The abstract definitions below distil exactly these three properties.

\begin{definition}[Reflexivity]
A relation $R$ on $A$ is \emph{reflexive} if
\[ \forall a \in A,\quad a \mathrel{R} a. \]
\emph{Read aloud:} ``for all $a$ in $A$, $a$ is related to $a$.''
\end{definition}

\begin{definition}[Symmetry]
A relation $R$ on $A$ is \emph{symmetric} if
\[ \forall a, b \in A,\quad a \mathrel{R} b \implies b \mathrel{R} a. \]
\emph{Read aloud:} ``for all $a$ and $b$ in $A$, if $a$ is related to $b$ then $b$ is related to $a$.''
\end{definition}

\begin{definition}[Transitivity]
A relation $R$ on $A$ is \emph{transitive} if
\[ \forall a, b, c \in A,\quad (a \mathrel{R} b \wedge b \mathrel{R} c) \implies a \mathrel{R} c. \]
\emph{Read aloud:} ``for all $a$, $b$, $c$ in $A$, if $a$ is related to $b$ and $b$ is related to $c$, then $a$ is related to $c$.''
\end{definition}

\begin{definition}[Antisymmetry]
A relation $R$ on $A$ is \emph{antisymmetric} if
\[ \forall a, b \in A,\quad (a \mathrel{R} b \wedge b \mathrel{R} a) \implies a = b. \]
\emph{Read aloud:} ``for all $a$ and $b$ in $A$, if $a$ is related to $b$ and $b$ is related to $a$, then $a$ equals $b$.''
\end{definition}

\begin{definition}[Equivalence relation]
A relation $\sim$ on $A$ is an \emph{equivalence relation} if it is reflexive, symmetric, and transitive. For each $a \in A$ the \emph{equivalence class} of $a$ is
\[ [a] \;=\; \{ x \in A : a \sim x \}, \]
\emph{Read aloud:} ``the class of $a$ equals the set of all $x$ in $A$ such that $a$ is equivalent to $x$.''
and the \emph{quotient set} is $A/{\sim} = \{ [a] : a \in A \}$.
\end{definition}

\begin{definition}[Partition]
A \emph{partition} of $A$ is a family $\mathcal{P} \subseteq \mathcal{P}(A)$ of nonempty subsets of $A$ such that
\begin{enumerate}[label=(P\arabic*)]
  \item $\bigcup_{B \in \mathcal{P}} B = A$ (covering);
  \item $B, B' \in \mathcal{P}$ with $B \neq B' \implies B \cap B' = \emptyset$ (pairwise disjoint).
\end{enumerate}
\end{definition}

\begin{definition}[Partial order]
A relation $\le$ on $A$ is a \emph{partial order} if it is reflexive, antisymmetric, and transitive. It is a \emph{total order} if additionally $\forall a, b \in A,\ a \le b \vee b \le a$.
\end{definition}

\section{Equivalence Relations Correspond to Partitions}

We now prove the central structural theorem of this concept.

\begin{lemma}\label{lem:classes-disjoint-or-equal}
Let $\sim$ be an equivalence relation on $A$, and let $a, b \in A$. Then $[a] = [b]$ if $a \sim b$, and $[a] \cap [b] = \emptyset$ otherwise.
\end{lemma}

\begin{proof}
Suppose $a \sim b$. If $x \in [a]$, then $a \sim x$, so by symmetry $x \sim a$, and combined with $a \sim b$ transitivity gives $x \sim b$, hence $b \sim x$ by symmetry, i.e. $x \in [b]$. The reverse inclusion is symmetric, so $[a] = [b]$.

Suppose instead $a \not\sim b$, but for contradiction take $x \in [a] \cap [b]$. Then $a \sim x$ and $b \sim x$. By symmetry $x \sim b$, and transitivity yields $a \sim b$, contradicting the hypothesis. Hence $[a] \cap [b] = \emptyset$.
\end{proof}

\begin{theorem}[Equivalence relations $\leftrightarrow$ partitions]\label{thm:eq-partition}
Let $A$ be a set. There is a bijection between the set of equivalence relations on $A$ and the set of partitions of $A$, given by
\[
\Phi(\sim) \;=\; \{ [a]_\sim : a \in A \},
\qquad
\Psi(\mathcal{P}) \;=\; \{ (a, b) \in A \times A : \exists B \in \mathcal{P},\ a, b \in B \}.
\]
\emph{Read aloud:} ``Phi of $\sim$ is the set of equivalence classes $[a]$ for $a$ in $A$; Psi of $\mathcal{P}$ is the set of pairs $(a,b)$ in $A \times A$ such that there exists a block $B$ in $\mathcal{P}$ containing both $a$ and $b$.''
\end{theorem}

\begin{proof}
\emph{Step 1: $\Phi(\sim)$ is a partition.}
For every $a \in A$, reflexivity gives $a \in [a]$, so each class is nonempty and $\bigcup_a [a] = A$. Pairwise disjointness of distinct classes is Lemma~\ref{lem:classes-disjoint-or-equal}.

\emph{Step 2: $\Psi(\mathcal{P})$ is an equivalence relation.}
Let $\approx := \Psi(\mathcal{P})$.
\begin{itemize}
  \item Reflexive: each $a \in A$ lies in some $B \in \mathcal{P}$ by (P1), so $a \approx a$.
  \item Symmetric: $a \approx b$ means $a, b \in B$ for some $B$; the same $B$ witnesses $b \approx a$.
  \item Transitive: if $a, b \in B$ and $b, c \in B'$, then $b \in B \cap B'$, forcing $B = B'$ by (P2); hence $a, c \in B$ and $a \approx c$.
\end{itemize}

\emph{Step 3: $\Phi \circ \Psi = \mathrm{id}$.}
Let $\mathcal{P}$ be a partition and set $\approx = \Psi(\mathcal{P})$. Pick any $a \in A$ and let $B_a \in \mathcal{P}$ be the (unique) block containing $a$. The class $[a]_\approx$ consists of all $x$ sharing some block with $a$, but uniqueness of the block makes this exactly $B_a$. Thus $\Phi(\Psi(\mathcal{P})) = \{ B_a : a \in A \} = \mathcal{P}$, since every block of $\mathcal{P}$ contains at least one element.

\emph{Step 4: $\Psi \circ \Phi = \mathrm{id}$.}
Let $\sim$ be an equivalence relation and set $\mathcal{P} = \Phi(\sim)$. For $a, b \in A$,
\[
a \,\Psi(\mathcal{P})\, b
\iff \exists [c] \in \mathcal{P},\ a, b \in [c]
\iff c \sim a \wedge c \sim b
\iff a \sim b,
\]
\emph{Read aloud:} ``$a$ is $\Psi(\mathcal{P})$-related to $b$ iff there exists a class $[c]$ in $\mathcal{P}$ containing both $a$ and $b$, iff $c$ is equivalent to $a$ and $c$ is equivalent to $b$, iff $a$ is equivalent to $b$.''
using symmetry and transitivity for the last equivalence. So $\Psi(\Phi(\sim)) = \sim$.

The two maps are mutually inverse, hence bijective.
\end{proof}

\section{An Example to Hold Onto}

\begin{example}[Congruence mod $n$]
Fix $n \in \mathbb{Z}_{>0}$. Define $a \equiv b \pmod n$ iff $n \mid (a - b)$. This is an equivalence relation on $\mathbb{Z}$, and Theorem~\ref{thm:eq-partition} shows the residue classes $[0], [1], \ldots, [n-1]$ form a partition of $\mathbb{Z}$ into exactly $n$ blocks.
\end{example}

\begin{example}[Divisibility]
On $\mathbb{N}_{>0}$, define $a \le_{\mid} b$ iff $a \mid b$. This is a partial order: reflexive ($a \mid a$), antisymmetric ($a \mid b \wedge b \mid a \Rightarrow a = b$ for positive integers), and transitive. It is \emph{not} total: $2 \nmid 3$ and $3 \nmid 2$.
\end{example}

\end{document}
